\chapter{Foundations of Generative Modeling}

\section{Introduction}

Generative modeling aims to learn the underlying distribution of data rather than just predicting labels. Instead of learning a mapping $x \mapsto y$, a generative model seeks to capture the structure of $p(x)$ or related distributions, enabling tasks such as sampling, density estimation, and representation learning.

This chapter establishes the conceptual and mathematical foundations for generative modeling. We introduce different types of generative models, clarify their objectives, and distinguish between likelihood-based, sampling-based, and representation-based perspectives.

\section{What It Means to Model a Data Distribution}

Let $X \sim p_{\text{data}}$ denote the unknown data-generating distribution. A generative model aims to learn a distribution $p_\theta(x)$ such that:
\[
p_\theta(x) \approx p_{\text{data}}(x).
\]

This can be approached in several ways:
\begin{itemize}
    \item estimating densities,
    \item generating samples,
    \item learning latent representations.
\end{itemize}

Each perspective emphasizes a different aspect of modeling.

\section{Joint, Marginal, and Conditional Generative Models}

Generative modeling can involve different distributions:

\subsection{Marginal Models}

Model $p(x)$ directly:
\[
p_\theta(x) \approx p_{\text{data}}(x).
\]

\subsection{Joint Models}

Model the joint distribution:
\[
p_\theta(x,y),
\]
which allows derivation of conditionals via:
\[
p(y|x) = \frac{p(x,y)}{p(x)}.
\]

\subsection{Conditional Models}

Model $p_\theta(x|y)$ or $p_\theta(y|x)$ directly. Examples include conditional generative models and sequence-to-sequence systems.

Each choice reflects a different modeling objective and application.

\section{Sampling, Likelihood, and Inference}

Generative models are often evaluated along three axes:

\subsection{Sampling}

Can we generate realistic samples $x \sim p_\theta(x)$?

\subsection{Likelihood}

Can we evaluate or approximate $p_\theta(x)$?

\subsection{Inference}

Can we infer latent structure, such as hidden variables $z$?

Different models prioritize different aspects:
\begin{itemize}
    \item likelihood-based models focus on density estimation,
    \item implicit models focus on sampling,
    \item latent-variable models emphasize representation.
\end{itemize}

\section{Latent-Variable Models}

Many generative models introduce latent variables $z$:
\[
p_\theta(x) = \int p_\theta(x|z)p(z)\,dz.
\]

\subsection{Interpretation}

\begin{itemize}
    \item $z$ captures underlying structure,
    \item $p(z)$ is a prior,
    \item $p_\theta(x|z)$ generates data from latent codes.
\end{itemize}

Latent-variable models enable:
\begin{itemize}
    \item structured generation,
    \item disentangled representations,
    \item dimensionality reduction.
\end{itemize}

\section{Explicit Density, Implicit Density, and Energy Views}

Generative models differ in how they represent distributions.

\subsection{Explicit Density Models}

These define $p_\theta(x)$ directly and allow likelihood evaluation.

Examples:
\begin{itemize}
    \item autoregressive models,
    \item normalizing flows,
    \item variational autoencoders (approximate likelihood).
\end{itemize}

\subsection{Implicit Models}

These define a sampling procedure:
\[
x = G_\theta(z), \quad z \sim p(z),
\]
without explicit density evaluation.

Example:
\begin{itemize}
    \item generative adversarial networks (GANs).
\end{itemize}

\subsection{Energy-Based Models}

These define:
\[
p_\theta(x) \propto e^{-E_\theta(x)},
\]
where $E_\theta(x)$ is an energy function.

They avoid explicit normalization but require specialized inference methods.

\section{Evaluation of Generative Models}

Evaluating generative models is challenging because there is no single metric.

\subsection{Likelihood-Based Evaluation}

\begin{itemize}
    \item log-likelihood,
    \item bits-per-dimension.
\end{itemize}

These measure how well the model fits the data distribution.

\subsection{Sample Quality}

\begin{itemize}
    \item visual inspection,
    \item metrics such as FID (in practice).
\end{itemize}

These measure realism of generated samples.

\subsection{Tradeoffs}

High likelihood does not always imply high-quality samples, and vice versa. This reflects differences between modeling density and capturing perceptual structure.

\section{Discriminative vs Generative Objectives}

Discriminative models learn:
\[
p(y|x),
\]
while generative models learn:
\[
p(x) \quad \text{or} \quad p(x,y).
\]

\subsection{Comparison}

\begin{itemize}
    \item discriminative models focus on prediction accuracy,
    \item generative models aim to capture data structure,
    \item generative models can simulate data and support unsupervised learning.
\end{itemize}

\section{A Unifying Perspective}

Generative modeling can be understood as learning a probability distribution over data. Different approaches emphasize:
\begin{itemize}
    \item likelihood (explicit density models),
    \item sampling (implicit models),
    \item structure (latent-variable models),
    \item energy (unnormalized models).
\end{itemize}

These perspectives are complementary and lead to different algorithmic approaches.

\section{Summary}

In this chapter we introduced the foundations of generative modeling. We defined the goal of modeling a data distribution and distinguished between marginal, joint, and conditional models. We discussed sampling, likelihood, and inference as key evaluation axes, and introduced latent-variable models.

We also classified generative models into explicit density, implicit, and energy-based approaches, and discussed challenges in evaluation. This framework sets the stage for specific generative models in the following chapters.